% ----------------------
\subsection{Section 81 (15 March 1832)}
% ----------------------
	\label{DC81}
	
	% -----
	\subsubsection{Historical Background}
	% -----
		\label{DC81-HB}
		
		% --
		\paragraph{JSP}
		% --
		
			``The index to Revelation Book 2 identifies `Jesse Gauze [Gause]' as the subject of this 15 March 1832 revelation, although the revelation refers to him only as `Jesse.' Gause had recently been called as a counselor to JS in the presidency of the high priesthood. A November 1831 revelation declared that presidents were to be appointed to preside over groups of men who held various offices, including the office of high priest. JS was ordained president of the high priesthood at a 25 January 1832 conference in Amherst, Ohio. On 8 March 1832, JS appointed Gause and Sidney Rigdon as `councillers of the ministry of the presidency of th[e] high Pristhood.' While Rigdon had been associated with JS almost since Rigdon's baptism in November 1830, Gause was a relatively new member, baptized sometime in late 1831 or early 1832. This 15 March 1832 revelation instructed Gause in his duties as counselor. In addition, it acknowledged JS as president of the high priesthood and stated that `the keys of the Kingdom' rested with him.
			
			``Because Gause apparently acted as a scribe for JS's revisions of the New Testament between 8 and 20 March 1832, it is likely that the nonextant original manuscript of this revelation was in Gause's own handwriting. Frederick G. Williams copied the revelation into Revelation Book 2, probably sometime before 1 April 1832. At a later point---sometime after the appointment of Williams as a counselor to JS in January 1833---Oliver Cowdery replaced Gause's name with Williams's in the Revelation Book 2 copy (which is the version featured below). Someone---likely Cowdery---crossed out Gause's name as well. The published versions of this revelation in the 1835 edition of the Doctrine and Covenants, the 1844 edition of the Doctrine and Covenants, and the 15 August 1844 issue of the \textit{Times and Seasons} all have Williams's name instead of Gause's. This indicates that JS and others regarded this revelation as containing general information about the duties of a counselor, rather than instructions specific to Gause.''
			
		% --
		\paragraph{HDDC}
		% --
		
			HDDC 1017--1021 has some background information on Gause and other aspects of this section; it's mainly because Woodford knew that Gause's name appeared crossed off in manuscript copies of the revelation, and he was trying to figure out why. There wasn't a lot of info available on Gause apparently.
		
	% -----
	\subsubsection{Versions}
	% -----
		\label{DC81-V}
		
		See the \href{https://drive.google.com/file/d/1goAvNz8jS4R8slYfMG_db55Ty2z_vmgK/view?usp=sharing}{sections comparison} document.
	
		\begin{compactdesc}
			\item[JSP] \hyperref[EMS-1832-JS-15-MAR-1832]{15 March 1832}
			\item[BC] ---
			\item[DC 1835] Section 79, 207
			\item[TS] 5:15 (15 August 1844), 609
			\item[DC 1844] Section 80, 317
			\item[MS] 14:10 (1 May 1852), 146
		\end{compactdesc}

	% -----
	\subsubsection{Section 81:1} % *DC81-1 
	% -----
		\label{DC81-1}

		VERILY, verily, I say unto you my servant Frederick G. Williams: Listen to the voice of him who speaketh, to the word of the Lord your God, and hearken to the calling wherewith you are called, even to be a high priest in my church, and a counselor unto my servant Joseph Smith, Jun.;

		% \begin{framed}
		% \scriptsize
			% \begin{compactdesc}
				% \item[JSP] 
				% % \item[BC] 
				% % \item[]
			% \end{compactdesc}
		% \end{framed}

	% -----
	\subsubsection{Section 81:2} % *DC81-2 
	% -----
		\label{DC81-2}

		Unto whom I have given the keys of the kingdom, which belong always unto the Presidency of the High Priesthood:

		% \begin{framed}
		% \scriptsize
			% \begin{compactdesc}
				% \item[JSP] 
				% % \item[BC] 
				% % \item[]
			% \end{compactdesc}
		% \end{framed}

	% -----
	\subsubsection{Section 81:3} % *DC81-3 
	% -----
		\label{DC81-3}

		Therefore, verily I acknowledge him and will bless him, and also thee, inasmuch as thou art faithful in counsel, in the office which I have appointed unto you, in prayer always, vocally and in thy heart, in public and in private, also in thy ministry in proclaiming the gospel in the land of the living, and among thy brethren.

		% \begin{framed}
		% \scriptsize
			% \begin{compactdesc}
				% \item[JSP] 
				% % \item[BC] 
				% % \item[]
			% \end{compactdesc}
		% \end{framed}

	% -----
	\subsubsection{Section 81:4} % *DC81-4 
	% -----
		\label{DC81-4}

		And in doing these things thou wilt do the greatest good unto thy fellow beings, and wilt promote the glory of him who is your Lord.

		% \begin{framed}
		% \scriptsize
			% \begin{compactdesc}
				% \item[JSP] 
				% % \item[BC] 
				% % \item[]
			% \end{compactdesc}
		% \end{framed}

	% -----
	\subsubsection{Section 81:5} % *DC81-5 
	% -----
		\label{DC81-5}

		Wherefore, be faithful; stand in the office which I have appointed unto you; succor the weak, lift up the hands which hang down, and strengthen the feeble knees.

		% \begin{framed}
		% \scriptsize
			% \begin{compactdesc}
				% \item[JSP] 
				% % \item[BC] 
				% % \item[]
			% \end{compactdesc}
		% \end{framed}

	% -----
	\subsubsection{Section 81:6} % *DC81-6 
	% -----
		\label{DC81-6}

		And if thou art faithful unto the end thou shalt have a crown of immortality, and eternal life in the mansions which I have prepared in the house of my Father.

		% \begin{framed}
		% \scriptsize
			% \begin{compactdesc}
				% \item[JSP] 
				% % \item[BC] 
				% % \item[]
			% \end{compactdesc}
		% \end{framed}

	% -----
	\subsubsection{Section 81:7} % *DC81-7 
	% -----
		\label{DC81-7}

		Behold, and lo, these are the words of Alpha and Omega, even Jesus Christ.  Amen.

		% \begin{framed}
		% \scriptsize
			% \begin{compactdesc}
				% \item[JSP] 
				% % \item[BC] 
				% % \item[]
			% \end{compactdesc}
		% \end{framed}